\section{Compact cylinder operators and a counting theorem}\label{sec:tree}
This section is independent of arithmetic. It shows that a sum of rank-one
operators attached to a tree of nested sets is compact, bounds its
eigenvalue count above and below by counts of vertices, and gives conditions
under which both bounds are asymptotic to the number of sets of mass above
the threshold. Section~\ref{sec:consequences} applies it to the tree of
Section~\ref{sec:kernel}.

Let $\cyl$ be a countable rooted tree, each vertex having finite depth.
On a probability space $(\Omega,\nu)$, associate to each vertex $v$ a
measurable set $A_v$ of positive mass $m_v=\nu(A_v)$, with $A_o=\Omega$ at
the root. Child sets lie inside the parent and are pairwise disjoint modulo
null sets. They need not cover the parent. Set
\[
 a_v=\sup_{c\text{ child of }v}\frac{m_c}{m_v},
 \qquad a_v=0\text{ at a leaf},\qquad
 D(t)=\#\{v:m_v>t\}.
\]
Throughout this section assume
\begin{equation}\label{eq:tree-assumptions}
 0\le a_v\le\theta<1,\qquad 0<\theta<1,\qquad D(t)<\infty\quad(t>0),
 \qquad M_\theta=\frac1{1-\theta}.
\end{equation}
For $f$ in a Hilbert space, $R_f$ denotes the positive rank-one operator
with quadratic form $\inner{x}{R_fx}=|\inner{x}{f}|^2$.

\begin{proposition}[Construction and compactness]\label{prop:compact}
The sum
\[
 T=\sum_{v\in\cyl}R_{\one_{A_v}}
\]
defines a positive compact operator on $L^2(\nu)$, with $\norm T\le M_\theta$.
The synthesis map $V:\ell^2(\cyl)\to L^2(\nu)$ is compact and satisfies
$T=VV^*$; the Gram operator $V^*V$ has entries $\nu(A_v\cap A_w)$.
If $T_v$ is the sum over descendants of $v$, including $v$, then
\[
 \norm{T_v}\le M_\theta m_v.
\]
For the finite truncation $T^{(\varepsilon)}=\sum_{m_v>\varepsilon}
R_{\one_{A_v}}$, with $0<\varepsilon<1$,
\begin{equation}\label{eq:cutoff-error}
 \norm{T-T^{(\varepsilon)}}\le M_\theta\varepsilon.
\end{equation}
\end{proposition}
\begin{proof}
At depth $\ell$, the sets are disjoint and their masses are at most
$\theta^\ell$. Thus the sum $Q_\ell$ at that depth is an orthogonal
sum of rank-one operators of norms $m_v$, and $\norm{Q_\ell}\le\theta^\ell$.
Each $Q_\ell$ is compact: for every positive threshold only finitely many
of these norms exceed it. Hence $\sum_{\ell\ge0}Q_\ell$ converges in operator
norm to a positive compact operator, with norm at most $M_\theta$.
The same argument within $A_v$ gives the subtree bound.

The family $\{v:m_v>\varepsilon\}$ is finite and ancestor closed.
The omitted vertices are the disjoint subtrees rooted at minimal omitted
vertices $c$, all of mass at most $\varepsilon$. Their operators act on
orthogonal spaces $L^2(A_c)$, proving \eqref{eq:cutoff-error}.

For any $f$, the sum of $|\inner{f}{\one_{A_v}}|^2$ is $\inner{f}{Tf}$.
This defines a bounded analysis operator $V^*$, whose adjoint is the stated
synthesis map. Then $VV^*=T$. Compactness of $T$ implies compactness of $V$,
for example by polar decomposition. The Gram entries follow on coordinate
vectors. These constructions also show that finite vertex sums converge
in norm: a finite set containing all vertices of mass above $\varepsilon$
leaves a positive remainder dominated by the one in \eqref{eq:cutoff-error}.
\end{proof}

Use the normalized measure $\nu/m_v$ on $A_v$ to define the normalized
subtree operator $\widehat T_v$. Extension after multiplication by
$m_v^{-1/2}$ gives a unitary $J_v$ onto $L^2(A_v,\nu)$, and
\[
 T_v=m_v J_v\widehat T_vJ_v^*.
\]
Decomposing into child spaces, with the unused part of $A_v$ assigned a
zero operator, gives
\begin{equation}\label{eq:normalized-decomposition}
 \widehat T_v=R_{\one}+
 \bigoplus_{c\text{ child of }v}\frac{m_c}{m_v}
 U_c\widehat T_cU_c^* .
\end{equation}
Here $U_c$ is the corresponding normalized embedding and
$\norm{\one}_{L^2(A_v,\nu/m_v)}=1$. Consequently
\begin{equation}\label{eq:subtree-norm}
 1\le\norm{\widehat T_v}\le1+M_\theta a_v\le M_\theta.
\end{equation}

\begin{proposition}[Tree sandwich]\label{prop:sandwich}
Let $\mathcal L_t$ be the leaves of the finite ancestor-closed tree
$\{v:m_v>t\}$. Then
\begin{equation}\label{eq:sandwich}
 |\mathcal L_t|\le\ncount(T;t)
 \le\#\{v:m_v\norm{\widehat T_v}>t\}.
\end{equation}
In particular,
\begin{equation}\label{eq:rough-sandwich}
 D(t)-D(t/\theta)\le\ncount(T;t)\le D((1-\theta)t).
\end{equation}
\end{proposition}
\begin{proof}
Leaf vertices are incomparable, so the compression of $V^*V$ to their
coordinates is diagonal with entries $m_v>t$. The min--max principle and
equality of the nonzero spectra of $V^*V$ and $VV^*$ prove the lower bound.

Write $n_v(u)=\ncount(\widehat T_v;u)$. Positive rank-one interlacing in
\eqref{eq:normalized-decomposition} gives
\[
 n_v(u)\le
 \one_{\{u<\norm{\widehat T_v}\}}
 +\sum_{c\text{ child of }v}n_c(um_v/m_c).
\]
For $u\ge\norm{\widehat T_v}$ both sides' child counts vanish, which explains
the indicator rather than an unrestricted $1$. For a fixed $u>0$, iteration
terminates uniformly: at depth $h$ the threshold is at least $u\theta^{-h}$,
and all normalized operator norms are at most $M_\theta$. Choose $h$ with
$u\theta^{-h}\ge M_\theta$. All contributing vertices have
$m_w>um_v/M_\theta$, so only finitely many can contribute.
Unrolling proves the upper bound at the root.

A nonleaf of $\{m_v>t\}$ has a child of mass above $t$, hence has mass
above $t/\theta$. This proves the first rough bound. The second follows
from \eqref{eq:subtree-norm}.
\end{proof}

\begin{definition}\label{def:regvar}
A positive function $D$ is regularly varying at zero with index $-\alpha$
if, for every $b>0$, $D(bt)/D(t)\to b^{-\alpha}$ as $t\downarrow0$.
We say $a_v\to0$ outside finite sets if
$\{v:a_v>\eta\}$ is finite for every $\eta>0$.
\end{definition}
The second condition is stronger than convergence of $a_v$ to zero along
each infinite path of the tree. It requires that, for each $\eta>0$, all
but finitely many vertices have every child of relative mass at most
$\eta$.

\begin{theorem}[Transfer from masses to eigenvalues]\label{thm:transfer}
Under \eqref{eq:tree-assumptions}, suppose $a_v\to0$ outside finite sets
and $D$ is regularly varying at zero with index $-\alpha$, where $\alpha>0$.
Then
\[
 \ncount(T;t)\sim D(t)\qquad(t\downarrow0).
\]
The conclusion holds as well for the nonzero spectrum of the Gram operator
$V^*V$.
\end{theorem}
\begin{proof}
Fix $0<\eta<1$ and let $F_\eta=\{v:a_v>\eta\}$.
Outside $F_\eta$, a nonleaf of $\{m_v>t\}$ has mass above $t/\eta$, and
$\norm{\widehat T_v}\le1+M_\theta\eta$. Proposition~\ref{prop:sandwich} gives
\[
 D(t)-|F_\eta|-D(t/\eta)
 \le\ncount(T;t)
 \le |F_\eta|+D\bigl(t/(1+M_\theta\eta)\bigr).
\]
Divide by $D(t)$ and let $t\downarrow0$. Since $\alpha>0$, $D(t)\to\infty$.
The lower limiting ratio is at least $1-\eta^\alpha$, and the upper
limiting ratio is at most $(1+M_\theta\eta)^\alpha$. Let $\eta\downarrow0$.
\end{proof}
